\mathtgt
\section{The Idea in Miniature}\label{sec:miniature}

This section develops the whole proof on a fragment with functions and type members only. The fragment has no fields, no recursive types, and no intersections. It is enough to express bad bounds and to make application the one rule of the machine that can get stuck. It is also small enough to state every definition. Section~\ref{sec:target} adds the rest of \FCdot one feature at a time, and Sections~\ref{sec:canonical} to~\ref{sec:safety} redo the proof, the translation, and the safety argument for it.

As Section~\ref{sec:intro} said, the fragment is not mechanized on its own. Section~\ref{sec:mechanization} lists where it departs from the mechanization.

\subsection{Overview}\label{sec:mini-overview}

Before the definitions, here is the bad-bounds example of Section~\ref{sec:problem} in \FCdot, step by step. The notation is defined in the subsections that follow. The point here is what each piece is for.

\paragraph{Types.}
A type member declaration becomes an object type: a list of propositions about the object, which the object type calls $\self$. The bad-bounds declaration $\dtyp{A}{\top}{\bot}$ becomes
\[
D = \obj{[\,\ple{\top}{\sel{\self}{A}},\ \ple{\sel{\self}{A}}{\bot}\,]} ,
\]
which says that the member $A$ of the object lies between $\top$ and $\bot$. Here $\mu$ only binds the self, and $\sel{\self}{A}$ is the member $A$ of the self. It is not a recursive type. Inside a type, $\ple{S}{T}$ is an assumption about the self, and evidence proves $S \le T$.

\paragraph{Evidence.}
Every use of subtyping is a proof term. In the body of $\src{\lambda(x : \dtyp{A}{\top}{\bot}).\ u}$, DOT derives $\src{\top \sub x.A \sub \bot}$. In \FCdot that derivation is the evidence
\[
\ev{trans}\ (\ev{member}\,x\,(\ev{refl}\,D)\,0)\ (\ev{member}\,x\,(\ev{refl}\,D)\,1)\ :\ \top \le \bot ,
\]
where $\ev{member}\,x\,(\ev{refl}\,D)\,i$ reads proposition $i$ of $D$, with $x$ for $\self$. Its middle argument is evidence that the type of $x$ is included in $D$, here $\ev{refl}\,D$. This evidence is well typed. Nothing rules it out, and nothing has to, because it is evidence under an assumption, the binder $x : D$. In System FC terms, $x : D$ is like a coercion variable $\run{\gamma : \mathit{Int} \sim \mathit{Bool}}$. The body may use it freely, because progress in FC only asks that no coercion in the empty context relates two distinct type constructors.

\paragraph{Stores.}
Now apply the lambda. Its argument is an atom, a variable under casts, rooted at a location $z$ of the store. Say the store holds at $z$ a literal that defines one member, $A$, as some type $T$. The store types $z$ at the literal's precise type $\obj{[\,\peq{\sel{\self}{A}}{T}\,]}$, which says only that $A$ is $T$. To be passed at the type $D$, the argument is the location under one or more casts. The inclusion evidence that a cast carries is called a coercion. The coercions of these casts compose to some $c : \obj{[\,\peq{\sel{\self}{A}}{T}\,]} \le D$.

\paragraph{Normal forms.}
The coercion $c$ knows nothing about $\sel{\self}{A}$ except what the precise type it starts from states, $\peq{\sel{\self}{A}}{T}$. So its proof of $\ple{\top}{\sel{\self}{A}}$ has to be a proof of $\top \le T$ followed by that fact. Its proof of $\ple{\sel{\self}{A}}{\bot}$ has to be that fact followed by a proof of $T \le \bot$. Over the store, evidence normalizes to a head form, which says how one type is included in another. Either the first is $\bot$, or the second is $\top$, or the two are equal after following definitions, or both are function types, or both are object types. Composing the two proofs would give a head form for $\top \le \bot$, and no such head form exists. So $c$ does not exist, the lambda is never applied, and its body, with its proof of $\top \le \bot$, never runs. Example~\ref{ex:mini-bad} gives the argument with every step justified.

The normalizer needs the store. In the body, $x$ is a parameter and nothing defines $\sel{x}{A}$, so the evidence there has no typed head form. Nothing normalizes it until an application replaces $x$ by an atom over the store (Section~\ref{sec:mini-theorem}).

\paragraph{A program that runs.}
Bad bounds are about a program that must never run. The section also follows one that does: the polymorphic identity $\mathit{id}$ of Section~\ref{sec:dot} applied to the literal $\tgt{\nu(\self.\,A = \mathit{Int})}$, which the program binds to $t$. As there, $\mathit{id}$ and $n : \mathit{Int}$ are variables bound earlier, and in rules $t$ still ranges over terms. The literal has the precise type $\mathit{Lit} = \obj{[\,\peq{\sel{\self}{A}}{\mathit{Int}}\,]}$. The let binds $t$ at $R$, the object type of the declaration $\dtyp{A}{\mathit{Int}}{\mathit{Int}}$, and $\mathit{id}$ takes a parameter of type $\mathit{Abs}$, the object type of $\dtyp{A}{\bot}{\top}$:
\[
R = \obj{[\,\ple{\mathit{Int}}{\sel{\self}{A}},\ \ple{\sel{\self}{A}}{\mathit{Int}}\,]} , \qquad
\mathit{Abs} = \obj{[\,\ple{\bot}{\sel{\self}{A}},\ \ple{\sel{\self}{A}}{\top}\,]} .
\]
The translation casts the literal from $\mathit{Lit}$ to $R$ and the argument of $\mathit{id}$ from $R$ to $\mathit{Abs}$, and it passes $n$ and the result through the two bounds of $t$. Section~\ref{sec:mini-evidence} gives the translated program.

\paragraph{Roadmap.}
Sections~\ref{sec:mini-types} to~\ref{sec:mini-machine} define the calculus: types, terms, evidence, and the machine. Sections~\ref{sec:mini-normalizer} and~\ref{sec:mini-typed} define the normalizer and say when its result is typed. The rest follows the textbook proof of Section~\ref{sec:textbook}. Its inversion lemma becomes Lemma~\ref{lem:mini-shapes}, which is about head forms and holds by inspection. Its \rn{Trans} case becomes Lemma~\ref{lem:mini-combine}, which composes head forms. Theorem~\ref{thm:mini-canon} in Section~\ref{sec:mini-theorem} normalizes evidence typed over a typed store to a typed head form. The canonical forms lemma becomes Corollary~\ref{cor:mini}. Section~\ref{sec:mini-safety} derives preservation and progress, and Section~\ref{sec:mini-translation} translates the DOT fragment.

\FloatBarrier
\subsection{Types: block names and telescopes}\label{sec:mini-types}

Two decisions shape the types of \FCdot. The first is how a type refers to a term. As in DOT, it does so only through a name $\sel{x}{A}$ for the type member $A$ of a variable $x$. We write $A$ and $B$ for labels, which is what DOT calls the names of type members. We call $\sel{x}{A}$ a block name. Every variable $x$ has one for every label, whether or not its type mentions the label, and together they form the block of $x$. \FCdot has no type variables, and block names take their place, as $\src{x.A}$ took the place of $\run{\alpha}$ in Section~\ref{sec:dot}. Binding $x$ brings its whole block into scope, as $\run{\Lambda\alpha}$ brings $\run{\alpha}$ into scope, and a binder may define its block or leave it abstract (Section~\ref{sec:mini-terms}). The self of an object type is a variable too, so $\sel{\self}{A}$ in the type $D$ of the overview is a block name of the self.

The name $\sel{x}{A}$ does not depend on the type of $x$. \FCdot retypes a variable by wrapping it in a cast (Section~\ref{sec:mini-terms}). A type never contains a cast, only the variable under it. So a cast changes what is known about $\sel{x}{A}$, never which name it is. For a lambda parameter or a let, what is known is what the type of $x$ says about $A$, as in DOT, and the evidence $\ev{member}$ of Section~\ref{sec:mini-evidence} reads it off. Example~\ref{ex:mini-cast} gives a derivation where the name matters. Section~\ref{sec:counterexample} shows that reading a declaration as an existential package, with $\src{x.A}$ the type variable obtained by opening it, cannot translate that derivation.

A second decision is how an object type says what is known. An object type is a telescope of propositions about a self binder:
\[
\begin{array}{@{}lrcl@{}}
\text{Types} & S,T & ::= & \bot \mid \sel{x}{A} \mid \pit[x]{S}{T} \mid \obj[x]{\tel} \\
\text{Propositions} & P & ::= & \ple{S}{T} \mid \peq{S}{T} \\
\text{Telescopes} & \tel & ::= & \tnil \mid \tcons{\tel}{P}
\end{array}
\]
Why propositions? Section~\ref{sec:problem} observed that a declaration is an assumption, in fact two, $\src{S \sub x.A}$ and $\src{x.A \sub T}$. An object type lists such assumptions one by one, so that evidence can cite each by its position. The elimination $\ev{member}$ reads one, and an object coercion uses them the same way (Section~\ref{sec:mini-evidence}). Why a self? A DOT declaration $\dtyp{A}{S}{T}$ is about the member $A$ of whatever object has it, without naming that object. The self names it, so that a proposition can say $\sel{\self}{A}$.

The object type $\obj[x]{\tel}$ binds its self $x$ in $\tel$, the way $\lambda(x : S).\,t$ binds $x$ in $t$. The propositions of $\tel$ mention it only through its block names $\sel{x}{A}$, like any other variable. The dependent function type $\pit[x]{S}{T}$ binds its parameter $x$ in $T$, and $T$ may mention the block of $x$.

An inclusion $\ple{S}{T}$ says that $S$ is included in $T$, and an equality $\peq{S}{T}$ that they are equal. They are propositions assumed about the self. The judgments $S \le T$ and $S \equiv T$ are what evidence proves (Section~\ref{sec:mini-evidence}). Propositions do not bind, so a telescope is a plain list $[\,P_0, \ldots, P_{n-1}\,]$. We write $\tat{\tel}{i}{P}$ when its $i$-th proposition is $P$, counting from $0$ at the left.

Bound variables can be renamed. In rules and examples the self is named $\self$ and the parameter of a function type $x$, unless a binder of that name is already in scope. A telescope written on its own, as evidence carries it, is the body of an object type, and its self is $\self$. The empty object type asserts nothing, so every type is included in it, and it serves as $\top$:
\[
\tgt{\top} \mathrel{:=} \obj{\tnil} .
\]
A type member declaration $\dtyp{A}{S}{T}$ corresponds to the object type with one proposition per bound,
\[
\obj{[\,\ple{S}{\sel{\self}{A}},\ \ple{\sel{\self}{A}}{T}\,]} .
\]
The source fragment has no recursive types, so its bounds do not mention the self. Formally their translations are weakened under the self binder, $\weak{S}$ and $\weak{T}$, and examples leave that weakening implicit. The types $D$, $R$, and $\mathit{Abs}$ of the overview are the cases of the bounds $\top..\bot$, $\mathit{Int}..\mathit{Int}$, and $\bot..\top$. As in DOT, nothing requires $S \le T$ of the bounds, so $D$ is a legal type. Section~\ref{sec:mini-translation} gives the translation of the source fragment.

\FloatBarrier
\subsection{Atoms, terms, and contexts}\label{sec:mini-terms}

\FCdot has no subsumption. Where DOT retypes a variable by subsumption, \FCdot wraps it in a cast. A variable under casts is an atom, and wherever DOT has a variable, \FCdot has an atom. Terms are in let-normal form over atoms:
\[
\begin{array}{@{}lrcl@{}}
\text{Atoms} & a, b & ::= & x \mid \ev{cast}\,a\,e \\
\text{Terms} & t, u & ::= & a \mid v \mid a\,b \mid \klet x = t \kin u \mid \ev{cast}\,t\,e \\
\text{Values} & v & ::= & \lambda(x : S).\,t \mid \tgt{\nu(x.\,W)} \mid \ev{cast}\,v\,e \\
\text{Witnesses} & W & ::= & \tnil \mid W, A = T \\
\text{Contexts} & \Gamma & ::= & \cdot \mid \Gamma, x : T \mid \Gamma, x : T \mathrel{:=} W
\end{array}
\]
The three casts are distinct constructors that are written alike. The cast of a term runs the term and then casts its result. The cast of a value and the cast of an atom are finished, like a lambda or a variable. The root of an atom is the variable underneath, $\rt{x} = x$ and $\rt{\ev{cast}\,a\,e} = \rt{a}$. A cast erases to nothing, so an atom is a variable at run time. Section~\ref{sec:mini-evidence} introduces the evidence $e$ of a cast.

A literal $\tgt{\nu(x.\,W)}$ in the fragment has only type members, given by its witness list $W$. It binds its self $x$ in $W$, as an object type does. A witness list defines every label it does not mention as $\top$. We write $W(A)$ for the witness at $A$. The precise type of a literal, $\obj{\tgt{\mathit{tel}(W)}}$, has one equality $\peq{\sel{\self}{A}}{T}$ for each entry $A = T$ of $W$, and nothing else. It says what the literal defines and no more.

A context has two kinds of binder. An opaque binder $x : T$ is what a lambda parameter or a let introduces, and its block is abstract. A transparent binder $x : T \mathrel{:=} W$ comes with a witness list $W$ that defines its block, and $\Gamma(x) \ni A \mathrel{:=} U$ says that the witness at $A$ is $U$. Transparent binders arise in the store, where every location is bound transparently by the value stored there.

For example, the literal $\tgt{\nu(\self.\,A = \mathit{Int})}$ has the witness list $[\,A = \mathit{Int}\,]$ and the precise type $\mathit{Lit}$. Stored at $t$, it gives the transparent binder $t : \mathit{Lit} \mathrel{:=} [\,A = \mathit{Int}\,]$. So $\sel{t}{A}$ is defined as $\mathit{Int}$, and every other name of $t$ as $\top$. The parameter $x$ of $\mathit{id}$ is opaque, and nothing defines $\sel{x}{A}$.

\paragraph{Resolution.}
Resolution follows definitions until it reaches a type that is not a defined name:
\[
\begin{array}{@{}lcl@{}}
\resolve{\Gamma}{\sel{x}{A}} & = & \resolve{\Gamma}{U} \quad \text{if } \Gamma(x) \ni A \mathrel{:=} U \\
\resolve{\Gamma}{T} & = & T \quad \text{otherwise}
\end{array}
\]
With $t$ bound as above, $\sel{t}{A}$ resolves to $\mathit{Int}$, and $\sel{t}{B}$ to $\top$. In the body of $\mathit{id}$, where $x$ is opaque, $\sel{x}{A}$ resolves to itself. Every name of a location is defined, so over a store resolution never stops at a name. Definitions may be cyclic, as in the literal $\tgt{\nu(\self.\,A = \sel{\self}{A})}$, which the typing of literals in Section~\ref{sec:mini-machine} allows. Section~\ref{sec:resolution} adds fuel, so that a cycle resolves to $\top$, which asserts nothing, and shows that both equations still hold.

\FloatBarrier
\subsection{Evidence}\label{sec:mini-evidence}

\begin{figure}[t]
\small
\color{tgtcol}
\begin{mathpar}
\inferrule{ }{\atomj{\Gamma}{x}{\Gamma(x)}}
\and
\inferrule{\atomj{\Gamma}{a}{S} \\ \lej{\Gamma}{e}{S}{T}}{\atomj{\Gamma}{\ev{cast}\,a\,e}{T}}
\\
\inferrule{ }{\lej{\Gamma}{\ev{refl}\,T}{T}{T}}
\and
\inferrule{\lej{\Gamma}{e}{S}{M} \\ \lej{\Gamma}{f}{M}{T}}{\lej{\Gamma}{\ev{trans}\,e\,f}{S}{T}}
\and
\inferrule{ }{\lej{\Gamma}{\ev{top}\,T}{T}{\top}}
\and
\inferrule{ }{\lej{\Gamma}{\ev{bot}\,T}{\bot}{T}}
\and
\inferrule{\eqj{\Gamma}{\varphi}{S}{T}}{\lej{\Gamma}{\ev{eqToLe}\,\varphi}{S}{T}}
\and
\inferrule{\lej{\Gamma}{e}{S_2}{S_1} \\ \lej{\Gamma, x : S_2}{f}{T_1}{T_2}}{\lej{\Gamma}{\ev{pi}\,e\,f}{\pit{S_1}{T_1}}{\pit{S_2}{T_2}}}
\and
\inferrule{\atomj{\Gamma}{a}{S} \\ \lej{\Gamma}{e}{S}{\obj{\tel}} \\ \tat{\tel}{i}{\ple{S'}{T'}}}
  {\lej{\Gamma}{\ev{member}\,a\,e\,i}{\inst{S'}{\rt{a}}}{\inst{T'}{\rt{a}}}}
\and
\inferrule{\morj{\Gamma}{m}{\tel}{\tel'}}{\lej{\Gamma}{\ev{obj}\,\tel\,m}{\obj{\tel}}{\obj{\tel'}}}
\\
\inferrule{ }{\eqj{\Gamma}{\ev{refl}\,T}{T}{T}}
\and
\inferrule{\eqj{\Gamma}{\varphi}{S}{T}}{\eqj{\Gamma}{\ev{symm}\,\varphi}{T}{S}}
\and
\inferrule{\eqj{\Gamma}{\varphi}{S}{M} \\ \eqj{\Gamma}{\psi}{M}{T}}{\eqj{\Gamma}{\ev{trans}\,\varphi\,\psi}{S}{T}}
\and
\inferrule{\Gamma(x) \ni A \mathrel{:=} T}{\eqj{\Gamma}{\ev{def}\,x\,A}{\sel{x}{A}}{T}}
\end{mathpar}
\color{black}
\caption{Typing of atoms and evidence in the fragment. The morphism judgment of the rule for $\ev{obj}$ is in Figure~\ref{fig:mini-templates}.}
\label{fig:mini-evidence}
\end{figure}

Evidence is a proof term for an inclusion or an equality. The fragment has
\[
\begin{array}{@{}lrcl@{}}
\text{Inclusion} & e & ::= & \ev{refl}\,T \mid \ev{trans}\,e\,e \mid \ev{top}\,T \mid \ev{bot}\,T \mid \ev{eqToLe}\,\varphi \mid \ev{pi}\,e\,e \mid \ev{member}\,a\,e\,i \mid \ev{obj}\,\tel\,m \\
\text{Equality} & \varphi & ::= & \ev{refl}\,T \mid \ev{symm}\,\varphi \mid \ev{trans}\,\varphi\,\varphi \mid \ev{def}\,x\,A
\end{array}
\]
where $\ev{obj}\,\tel\,m$ is an object coercion, explained below. Inclusion and equality evidence are separate syntax, and $\ev{refl}$ and $\ev{trans}$ name a constructor of each. Figure~\ref{fig:mini-evidence} gives the typing. It includes the typing of atoms, which is mutual with that of evidence, because a cast atom carries evidence and $\ev{member}$ takes an atom. An inclusion judgment $\lej{\Gamma}{e}{S}{T}$ says that $e$ proves $S \le T$ in $\Gamma$, and an equality judgment $\eqj{\Gamma}{\varphi}{S}{T}$ says that $\varphi$ proves $S \equiv T$. The endpoints are determined by the evidence and the context, so evidence can be checked without inference, by a checker in the spirit of Core Lint. Equality evidence enters inclusion evidence only through $\ev{eqToLe}$.

Most rules are what one expects. $\ev{refl}$ and $\ev{trans}$ translate \rn{Refl} and \rn{Trans} and make inclusion a preorder. $\ev{top}$ and $\ev{bot}$ translate \rn{Top} and \rn{Bot} and give its extremes. $\ev{pi}$ translates \rn{All}, contravariant in the domain, with the codomains compared under a parameter of the smaller domain. Unlike equality evidence, inclusion evidence has no symmetry rule. Three rules need explanation: $\ev{member}$, $\ev{obj}$, and $\ev{def}$.

\paragraph{Elimination at an atom.}
The rule $\ev{member}$ is the one rule that connects the type of a term to the block of that term. It translates the selection rules \rn{<:-Sel} and \rn{Sel-<:}. Given an atom $a$ of type $S$ and evidence $e$ that $S$ is included in an object type $\obj{\tel}$, the $i$-th proposition of $\tel$ holds of the root of $a$. The proposition mentions the self, and the rule instantiates the self by the root, written $\inst{S'}{\rt{a}}$. With $e = \ev{refl}$ the rule reads the declared type of a variable, as the overview did with $\ev{member}\,x\,(\ev{refl}\,D)\,0 : \top \le \sel{x}{A}$.

The rule takes an atom, not just a variable, because the translation of a variable typing is an atom (Section~\ref{sec:mini-translation}) and the machine substitutes atoms for variables (Section~\ref{sec:mini-machine}). Once $\klet y = \ev{cast}\,z\,c \kin u$ has run, the evidence $\ev{member}\,y\,e\,i$ in its body has become $\ev{member}\,(\ev{cast}\,z\,c)\,e\,i$. The translation always passes $\ev{refl}$ for $e$.

\begin{example}[A block name under a cast]\label{ex:mini-cast}
This example shows that a cast changes what is known about a name, never which name it is. Take the type $\mathit{Abs}$ of the overview, and let
\[
\begin{array}{@{}lcl@{}}
\mathit{Abs}_{\mathit{Int}} & = & \obj{[\,\ple{\mathit{Int}}{\sel{\self}{A}},\ \ple{\sel{\self}{A}}{\top}\,]} \\
\mathit{Box} & = & \obj{[\,\ple{\mathit{Abs}}{\sel{\self}{B}},\ \ple{\sel{\self}{B}}{\mathit{Abs}_{\mathit{Int}}}\,]}
\end{array}
\]
and $\Gamma = x : \mathit{Box},\ w : \mathit{Abs}$. The evidence
\[
e_x = \ev{trans}\ (\ev{member}\,x\,(\ev{refl}\,\mathit{Box})\,0)\ (\ev{member}\,x\,(\ev{refl}\,\mathit{Box})\,1) \ :\ \mathit{Abs} \le \mathit{Abs}_{\mathit{Int}}
\]
goes through the two bounds of $x$, so the atom $\ev{cast}\,w\,e_x$ has type $\mathit{Abs}_{\mathit{Int}}$. Its root is $w$, and eliminating the lower bound of $\mathit{Abs}_{\mathit{Int}}$ there gives
\[
\ev{member}\ (\ev{cast}\,w\,e_x)\ (\ev{refl}\,\mathit{Abs}_{\mathit{Int}})\ 0 \ :\ \mathit{Int} \le \sel{w}{A} .
\]
In DOT this is the derivation of $\src{\mathit{Int} \sub w.A}$ from $\src{w : \dtyp{A}{\mathit{Int}}{\top}}$, which subsumption gets from $\src{w : \dtyp{A}{\bot}{\top}}$ and the two bounds of $\src{x.B}$. The result is about $\sel{w}{A}$, the name that the declared type of $w$ talks about, although the bound $\mathit{Int}$ was learned through $\mathit{Abs}_{\mathit{Int}}$. The evidence $\ev{member}\,w\,e_x\,0$ has the same type, so what matters is the name, not the atom. Like $D$ in the overview, $\mathit{Box}$ has bad bounds, which a binder may assume.
\end{example}

\paragraph{Coercions between object types.}
An object type is a list of propositions about its self. A coercion from $\obj{\tel}$ to $\obj{\tel'}$ has to prove every proposition of $\tel'$ about the same self, and all it may use about the self is what $\tel$ says. For an inclusion $S \le T$ we call $S$ its source and $T$ its target, and $\tel$ and $\tel'$ are the source and target telescopes of the coercion. Here is the typical case, the counterpart of DOT's rule \rn{Typ} (Figure~\ref{fig:source-typing}), contravariant in the lower bound and covariant in the upper. Let
\[
\tel = [\,\ple{S_1}{\sel{\self}{A}},\ \ple{\sel{\self}{A}}{T_1}\,] , \qquad \tel' = [\,\ple{S_2}{\sel{\self}{A}},\ \ple{\sel{\self}{A}}{T_2}\,] ,
\]
with evidence $d_1 : S_2 \le S_1$ and $d_2 : T_1 \le T_2$ about types that do not mention the self. Each proposition of $\tel'$ follows from one proposition of $\tel$ by stretching one of its ends:
\[
\begin{array}{@{}l@{\qquad}l@{}}
S_2 \le S_1 \sqsubseteq \sel{\self}{A} & \text{by } d_1 \text{, then proposition } 0 \text{ of } \tel \\
\sel{\self}{A} \sqsubseteq T_1 \le T_2 & \text{by proposition } 1 \text{ of } \tel \text{, then } d_2
\end{array}
\]
\begin{figure}[t]
\small
\color{tgtcol}
\begin{mathpar}
\inferrule{\morj{\Gamma}{\tau_i}{\tel}{P_i} \text{ for each } i < n}{\morj{\Gamma}{[\,\tau_0, \ldots, \tau_{n-1}\,]}{\tel}{[\,P_0, \ldots, P_{n-1}\,]}}
\and
\inferrule{\holej{\tel}{H}{X}{Y} \\ \Gamma \vdash q : S \le X \\ \Gamma \vdash q' : Y \le T}
  {\morj{\Gamma}{\ev{le}\,q\,H\,q'}{\tel}{\ple{S}{T}}}
\\
\inferrule{\tat{\tel}{j}{\ple{X}{Y}}}{\holej{\tel}{\ev{le}\,j}{X}{Y}}
\and
\inferrule{\tat{\tel}{j}{\peq{X}{Y}}}{\holej{\tel}{\ev{eq}\,j}{X}{Y}}
\and
\inferrule{\tat{\tel}{j}{\peq{Y}{X}}}{\holej{\tel}{\ev{eqSym}\,j}{X}{Y}}
\\
\inferrule{ }{\Gamma \vdash \ev{none} : X \le X}
\and
\inferrule{\lej{\Gamma}{e}{U_1}{U_2}}{\Gamma \vdash \ev{some}\,e : \weak{U_1} \le \weak{U_2}}
\end{mathpar}
\color{black}
\caption{Typing of morphisms, templates, holes, and sides. A morphism proves a telescope, and a template proves one proposition $P$, from the source telescope $\tel$. A hole $H$ reads the proposition of the source telescope that it names as an inclusion, $\holej{\tel}{H}{X}{Y}$. A side is typed in $\Gamma$, without the self binder, and its endpoints are weakened under it.}
\label{fig:mini-templates}
\end{figure}%
A template records such a proof. In general a template proves an inclusion $\ple{S}{T}$ of the target telescope in three links:
\[
S \;\underset{q}{\le}\; X \;\underset{H}{\sqsubseteq}\; Y \;\underset{q'}{\le}\; T .
\]
The middle link is a proposition of the source. Each outer link is either empty or evidence about types outside the object. The syntax of what was just drawn is
\[
\begin{array}{@{}lrcl@{}}
\text{Morphisms} & m & ::= & [\,\tau, \ldots, \tau\,] \\
\text{Templates} & \tau & ::= & \ev{le}\,q\,H\,q' \\
\text{Holes} & H & ::= & \ev{le}\,j \mid \ev{eq}\,j \mid \ev{eqSym}\,j \\
\text{Sides} & q & ::= & \ev{none} \mid \ev{some}\,e
\end{array}
\]
The template $\ev{le}\,q\,H\,q'$ names the three links, and $\ev{le}$ marks an inclusion, in templates and holes alike.
\begin{itemize}
\item The hole $H$ is the middle link. It names a proposition of the source by its position $j$ and says how to read it. $\ev{le}\,j$ reads an inclusion $\ple{X}{Y}$, $\ev{eq}\,j$ reads an equality $\peq{X}{Y}$ forwards as $\ple{X}{Y}$, and $\ev{eqSym}\,j$ reads it backwards as $\ple{Y}{X}$. It is called a hole because the normalizer later fills it with what is known about that proposition at a particular object (Section~\ref{sec:mini-normalizer}).
\item The sides $q$ and $q'$ are the outer links. A side is $\ev{none}$ when its link is empty, $S = X$ or $Y = T$, and $\ev{some}\,e$ for evidence $e : S \le X$ or $e : Y \le T$.
\end{itemize}
A morphism $m$ is the list of templates, one per proposition of the target, in order. A morphism from $\tel$ to $\tel'$ proves every proposition of $\tel'$ from $\tel$, as a substitution maps one context to another. The coercion $\ev{obj}\,\tel\,m$ carries its source telescope, so that its source can be read off the evidence. Figure~\ref{fig:mini-templates} gives the typing. For the two telescopes above,
\[
\ev{obj}\,\tel\,[\,\ev{le}\,(\ev{some}\,d_1)\,(\ev{le}\,0)\,\ev{none},\ \ev{le}\,\ev{none}\,(\ev{le}\,1)\,(\ev{some}\,d_2)\,] \ :\ \obj{\tel} \le \obj{\tel'} .
\]
For the running example, passing an object of type $R$ at the type $\mathit{Abs}$, to instantiate $\mathit{id}$ at $t$, needs the coercion
\[
\begin{array}{@{}l@{}}
\mathit{inst} = \ev{obj}\,[\,\ple{\mathit{Int}}{\sel{\self}{A}},\ \ple{\sel{\self}{A}}{\mathit{Int}}\,] \\
\qquad [\,\ev{le}\,(\ev{some}\,(\ev{bot}\,\mathit{Int}))\,(\ev{le}\,0)\,\ev{none},\ \ev{le}\,\ev{none}\,(\ev{le}\,1)\,(\ev{some}\,(\ev{top}\,\mathit{Int}))\,] \ :\ R \le \mathit{Abs} .
\end{array}
\]
Its two templates link $\bot \le \mathit{Int} \sqsubseteq \sel{\self}{A} \le \sel{\self}{A}$ and $\sel{\self}{A} \le \sel{\self}{A} \sqsubseteq \mathit{Int} \le \top$.

A present side $\ev{some}\,e$ is typed outside the object type and weakened under the self, so it is self-free: its endpoints do not mention the self. An absent side is the identity, and it may sit between types that mention the self, as $\sel{\self}{A} \le \sel{\self}{A}$ does in $\mathit{inst}$. So a template never proves anything about the self by evidence of its own. All it knows about the self is the one source proposition in its hole, and it only stretches that proposition at its ends. This restriction is what keeps normalization structural (Section~\ref{sec:mini-theorem}). In the fragment it loses nothing, since \rn{Typ} compares bounds and bounds do not mention the self. Section~\ref{sec:why-structural} describes an earlier design without it.

Holes that read an equality matter for literals. The precise type of a literal that defines $A$ as $T$ is $\obj{[\,\peq{\sel{\self}{A}}{T}\,]}$. The coercion to the declaration with the exact bounds $T$ reads both bounds off the one equality:
\[
\mathit{litCo}_T = \ev{obj}\,[\,\peq{\sel{\self}{A}}{T}\,]\,[\,\ev{le}\,\ev{none}\,(\ev{eqSym}\,0)\,\ev{none},\ \ev{le}\,\ev{none}\,(\ev{eq}\,0)\,\ev{none}\,] .
\]
It proves $\obj{[\,\peq{\sel{\self}{A}}{T}\,]} \le \obj{[\,\ple{T}{\sel{\self}{A}},\ \ple{\sel{\self}{A}}{T}\,]}$. The first template reads the equality backwards as $T \le \sel{\self}{A}$, the lower bound, and the second reads it forwards as the upper bound. Neither needs a side. The running example uses $\mathit{litCo} = \mathit{litCo}_{\mathit{Int}}$, from $\mathit{Lit}$ to $R$. With $\mathit{litCo}$ and $\mathit{inst}$, the program of the overview translates to
\[
\klet t = \ev{cast}\,\tgt{\nu(\self.\,A = \mathit{Int})}\,\mathit{litCo} \kin \klet g = \mathit{id}\,(\ev{cast}\,t\,\mathit{inst}) \kin \ev{cast}\,(g\,(\ev{cast}\,n\,\mathit{lo}))\,\mathit{hi} ,
\]
where $\mathit{lo} = \ev{member}\,t\,(\ev{refl}\,R)\,0$ and $\mathit{hi} = \ev{member}\,t\,(\ev{refl}\,R)\,1$ read the lower and the upper bound of $t$. Section~\ref{sec:mini-translation} derives it from a typing derivation.

\paragraph{Definitions.}
The rule $\ev{def}$ reads the definition of a name off a transparent binder. No program of the fragment contains it, because programs are typed before they run, under opaque binders. A running program reads what a literal defines through a cast like $\mathit{litCo}$ (Example~\ref{ex:mini-sel}). The fragment keeps $\ev{def}$ as the evidence counterpart of resolution, and the full calculus uses it to type the fields of a literal (Section~\ref{sec:translation}). In System FC terms, $\ev{def}$ is an axiom, and the store is the set of axioms.

\FloatBarrier
\subsection{Typing terms, stores, and the machine}\label{sec:mini-machine}

\begin{figure}[t]
\small
\color{tgtcol}
\[
\begin{array}{@{}lcl@{}}
\mst{\sigma}{K}{\klet x = t \kin u} &\step& \mst{\sigma}{K \triangleright \ev{let}\,u}{t} \\
\mst{\sigma}{K}{\ev{cast}\,t\,e} &\step& \mst{\sigma}{K \triangleright \ev{cast}\,e}{t} \\
\mst{\sigma}{K \triangleright \ev{cast}\,e}{v} &\step& \mst{\sigma}{K}{\ev{cast}\,v\,e} \\
\mst{\sigma}{K \triangleright \ev{cast}\,e}{a} &\step& \mst{\sigma}{K}{\ev{cast}\,a\,e} \\
\mst{\sigma}{K \triangleright \ev{let}\,u}{v} &\step& \mst{(\sigma, \mathit{core}(v))}{\weak{K}}{\mathit{adjust}(u, v)} \\
\mst{\sigma}{K \triangleright \ev{let}\,u}{a} &\step& \mst{\sigma}{K}{\inst{u}{a}} \\
\mst{\sigma}{K}{x\,b} &\step& \mst{\sigma}{K}{\inst{t_0}{b}} \\
& & \quad \text{if } \sigma(x) = \lambda(z : S_0).\,t_0 \\
\mst{\sigma}{K}{a\,b} &\step& \mst{\sigma}{K}{\inst{t_0}{b}} \\
& & \quad \text{if } \sigma(\rt{a}) = \lambda(z : S_0).\,t_0 \text{ and } \mathit{chain}_\sigma(a) = \form{id} \\
\mst{\sigma}{K}{a\,b} &\step& \mst{\sigma}{K}{\ev{cast}\,(\inst{t_0}{\ev{cast}\,b\,d})\,(\inst{c}{b})} \\
& & \quad \text{if } \sigma(\rt{a}) = \lambda(z : S_0).\,t_0 \text{ and } \mathit{chain}_\sigma(a) = \form{pi}\,d\,c
\end{array}
\]
\[
\begin{array}{@{}lcl@{}}
\mathit{core}(\ev{cast}\,v\,e) & = & \mathit{core}(v) \\
\mathit{core}(v) & = & v \quad \text{if } v \text{ is not a cast} \\
\mathit{adjust}(u, v) & = & u \quad \text{if } v \text{ is not a cast} \\
\mathit{adjust}(u, \ev{cast}\,(\cdots(\ev{cast}\,v_0\,e_1)\cdots)\,e_n) & = & u[x \mapsto \ev{cast}\,x\,\weak{\bar{e}}]
\end{array}
\]
\begin{mathpar}
\inferrule{ }{\contj{\Gamma}{\cdot}{T}{T}}
\and
\inferrule{\tmj{\Gamma, x : T}{u}{\weak{U}} \\ \contj{\Gamma}{K}{U}{U'}}{\contj{\Gamma}{K \triangleright \ev{let}\,u}{T}{U'}}
\and
\inferrule{\lej{\Gamma}{e}{T}{U} \\ \contj{\Gamma}{K}{U}{U'}}{\contj{\Gamma}{K \triangleright \ev{cast}\,e}{T}{U'}}
\and
\inferrule{\storej{\sigma}{\Gamma} \\ \tmj{\Gamma}{t}{T} \\ \contj{\Gamma}{K}{T}{U}}{\vdash \mst{\sigma}{K}{t} : U}
\end{mathpar}
\color{black}
\caption{The machine of the fragment. The second step is for the cast of a term, not for the cast of a value or of an atom (Section~\ref{sec:mini-terms}). The last two application steps apply when $a$ is not a variable, and Section~\ref{sec:mini-normalizer} defines $\mathit{chain}_\sigma(a)$ and the head forms $\form{id}$ and $\form{pi}\,d\,c$. In $\mathit{adjust}$ the value $v_0$ is not a cast and $\bar{e} = \ev{trans}\,(\cdots(\ev{trans}\,e_1\,e_2)\cdots)\,e_n$ is the composite of the casts. The body $u$ of a let frame binds the new variable $x$, and $\weak{K}$ is $K$ weakened under it. $\contj{\Gamma}{K}{T}{U}$ says that $K$ accepts a value of type $T$ and produces one of type $U$. A state is final when its continuation is empty and its term is a value or an atom.}
\label{fig:mini-machine}
\end{figure}

The typing of terms and values is
\begin{mathpar}
\color{tgtcol}
\inferrule{\atomj{\Gamma}{a}{T}}{\tmj{\Gamma}{a}{T}}
\and
\inferrule{\valj{\Gamma}{v}{T}}{\tmj{\Gamma}{v}{T}}
\and
\inferrule{\atomj{\Gamma}{a}{\pit{S}{T}} \\ \atomj{\Gamma}{b}{S}}{\tmj{\Gamma}{a\,b}{\inst{T}{\rt{b}}}}
\and
\inferrule{\tmj{\Gamma}{t}{T} \\ \tmj{\Gamma, x : T}{u}{\weak{U}}}{\tmj{\Gamma}{\klet x = t \kin u}{U}}
\and
\inferrule{\tmj{\Gamma}{t}{S} \\ \lej{\Gamma}{e}{S}{T}}{\tmj{\Gamma}{\ev{cast}\,t\,e}{T}}
\end{mathpar}
\begin{mathpar}
\color{tgtcol}
\inferrule{\tmj{\Gamma, x : S}{t}{T}}{\valj{\Gamma}{\lambda(x : S).\,t}{\pit{S}{T}}}
\and
\inferrule{ }{\valj{\Gamma}{\tgt{\nu(\self.\,W)}}{\obj{\tgt{\mathit{tel}(W)}}}}
\and
\inferrule{\valj{\Gamma}{v}{S} \\ \lej{\Gamma}{e}{S}{T}}{\valj{\Gamma}{\ev{cast}\,v\,e}{T}}
\end{mathpar}
Application instantiates the codomain at the root of the argument, as the source rule \rn{All-E} does. A let requires the type of the body not to mention the bound variable, as in the source. A literal has its precise type, and a literal used at another type is a literal under a cast.

\paragraph{Stores.}
A store holds one value per location. Store typing $\storej{\sigma}{\Gamma}$ types each entry in the context of the entries before it. Every entry is a lambda or a literal, not a cast. It is typed at its precise type, which for a lambda is its function type. Its location is bound transparently. The self of a stored literal becomes the new location, so its witnesses define the block of that location as they stand:
\begin{mathpar}
\color{tgtcol}
\inferrule{ }{\storej{\cdot}{\cdot}}
\and
\inferrule{\storej{\sigma}{\Gamma} \\ v \text{ is not a cast} \\ \valj{\Gamma}{v}{T}}{\storej{\sigma, v}{\Gamma, x : T \mathrel{:=} \mathit{witnesses}(v)}}
\end{mathpar}
Here $\mathit{witnesses}(\tgt{\nu(\self.\,W)}) = W$ and $\mathit{witnesses}(\lambda(x : S).\,t) = \tnil$, so every block name of a lambda is defined as $\top$.

The store is inert: a location's type is the precise type of its value and asserts nothing beyond its definitions. This is how the store discharges the assumptions of Section~\ref{sec:idea}. A definition holds by construction. A fact that might be false, such as $\top \le \sel{z}{A}$, never enters the store. A program can claim it only through a cast, as $c$ was in the overview, and Theorem~\ref{thm:mini-canon} normalizes the evidence of that cast against the definitions. Section~\ref{sec:comparison} relates this store to the inert contexts of Rapoport et al.

\paragraph{The machine.}
A state $\mst{\sigma}{K}{t}$ has a store, a continuation of let frames and cast frames, and a term. Figure~\ref{fig:mini-machine} gives all the steps, the typing of continuations and states, and the two functions that allocation uses. Most steps push and pop frames. Substitution of an atom for a variable, written $\inst{u}{a}$, or $u[x \mapsto a]$ when the bound variable matters, puts the atom wherever an atom is expected and only its root into types and into $\ev{def}$. Allocation keeps the store inert. It stores the value underneath the casts, $\mathit{core}(v)$, at its precise type, and the body uses the new variable through the composite of the stripped casts, $\mathit{adjust}(u, v)$.

In the running example, the let that binds $t$ runs its cast, and the literal comes back through the cast frame as the cast value $\ev{cast}\,\tgt{\nu(\self.\,A = \mathit{Int})}\,\mathit{litCo}$. At the let frame the machine stores the bare literal at $t$, at its precise type $\mathit{Lit}$, and continues with $u[t \mapsto \ev{cast}\,t\,\mathit{litCo}]$, where $u$ is the body of the let. The body was typed with $t : R$. It now sees $t$ through an atom of type $R$, and no type changes, because substitution puts only the root $t$ into types. The argument of $\mathit{id}$ becomes $b = \ev{cast}\,(\ev{cast}\,t\,\mathit{litCo})\,\mathit{inst}$, the location under two casts.

\paragraph{Application.}
The steps that matter are the last three of Figure~\ref{fig:mini-machine}. Application through a bare variable is $\beta$. Application through a cast atom is where evidence meets execution. The machine normalizes the chain of casts on the atom, $\mathit{chain}_\sigma(a)$, which Section~\ref{sec:mini-normalizer} defines. If the chain is $\form{id}$, the stored lambda is applied directly. If it is a function form $\form{pi}\,d\,c$, the argument is cast by the domain evidence $d$ and the result by the codomain evidence $c$. This is how a coercion between function types acts on a function. The coercion may reach the function type through selections and definitions, so the machine finds its parts by normalizing the chain at run time. All three steps erase to $\beta$ on the stored lambda. The casts are there because \FCdot has no subsumption, and the next state has to be typed as it stands. Progress therefore has to show that the normalizer succeeds and returns one of these two forms whenever an atom of function type is applied.

\FloatBarrier
\subsection{Head forms and the normalizer}\label{sec:mini-normalizer}

At an application the machine asks one question: what does the chain of casts on an atom normalize to? This subsection defines the normalizer that answers it. Over a store, every variable is a location, every location is bound transparently, and evidence can be normalized. The normal forms are head forms, entries, and views:
\[
\begin{array}{@{}lrcl@{}}
\text{Head forms} & F, G & ::= & \form{bot} \mid \form{top} \mid \form{id} \mid \form{pi}\,d\,c \mid \form{obj}\,\mathit{Es} \\
\text{Entries} & E & ::= & \form{le}\,F\,H\,G \\
\text{Entry lists} & \mathit{Es} & ::= & \tnil \mid \tcons{\mathit{Es}}{E} \\
\text{Proposition forms} & Q & ::= & \form{le}\,F \mid \form{eq} \\
\text{Views} & V & ::= & \tnil \mid \tcons{V}{Q}
\end{array}
\]
We write $\mathit{hnf}(e)$ for the head normal form of $e$. A hat marks a normal form, and tells the form $\form{obj}$ from the evidence $\ev{obj}$, and an entry $\form{le}$ from a template $\ev{le}$. $\form{bot}$ and $\form{top}$ say that the source is $\bot$ or the target is $\top$. $\form{id}$ says that the two endpoints are equal after following definitions. $\form{pi}\,d\,c$ is a function form. Its domain evidence $d$ and codomain evidence $c$ stay evidence. The machine uses them only as casts, and $c$ lives under the parameter, where there is no store to normalize against. Normalization stops at the binder, as weak head reduction stops at a $\lambda$. $\form{obj}\,\mathit{Es}$ is an object form, with one entry per target proposition. An entry $\form{le}\,F\,H\,G$ is a normalized template. Its sides have become head forms, with $\form{id}$ for an absent side, and its hole still names a source proposition.

\paragraph{Views.}
Head forms alone are not enough. A form for $S \le \obj{\tel}$ is relative. Its holes name propositions of $S$, which are still assumptions about whatever object has type $S$. To normalize $\ev{member}\,a\,e\,i$, the normalizer must know how those propositions hold for the object at the root of $a$. A view records this, with one proposition form per proposition of the atom's type: $\form{le}\,F$ for an inclusion and $\form{eq}$ for an equality. The view of a location is $\form{eq}$ at every position, since the precise type of its literal has only equalities and each holds by definition. Each cast transforms the view by the head form of its evidence. Through an object form, each entry looks up the position its hole names and wraps the result in its sides (Example~\ref{ex:mini-view}). So the view answers $\ev{member}$, and the chain answers application.

\paragraph{Composition.}
Composition $F \circ G$ is written in diagrammatic order, like $\ev{trans}$: if $F$ proves $S \le M$ and $G$ proves $M \le T$, then $F \circ G$ proves $S \le T$. It is defined by the first equation that applies:
\[
\begin{array}{@{}lcl@{}}
\form{id} \circ G & = & G \\
F \circ \form{id} & = & F \\
\form{bot} \circ G & = & \form{bot} \\
F \circ \form{top} & = & \form{top} \\
\form{pi}\,d_1\,c_1 \circ \form{pi}\,d_2\,c_2 & = & \form{pi}\,(\ev{trans}\,d_2\,d_1)\,(\ev{trans}\,(c_1[x \mapsto \ev{cast}\,x\,\weak{d_2}])\,c_2) \\
\form{obj}\,\mathit{Es}_1 \circ \form{obj}\,\mathit{Es}_2 & = & \form{obj}\,(\mathit{Es}_1 \text{ substituted into } \mathit{Es}_2) \\
F \circ G & = & F \quad \text{otherwise}
\end{array}
\]
Two function forms compose contravariantly in the domain. The first codomain evidence lives under a parameter of the middle domain. Substituting the parameter cast through $d_2$, weakened under the parameter, moves that evidence under a parameter of the final domain. Two object forms compose by substitution. Each entry $\form{le}\,F_2\,H_2\,G_2$ of $\mathit{Es}_2$ names a proposition of the middle telescope, and the entry of $\mathit{Es}_1$ at that position proves it from the source. If $H_2$ is $\ev{le}\,j$ and the entry of $\mathit{Es}_1$ at $j$ is $\form{le}\,F_1\,H_1\,G_1$, the substitution gives
\[
\form{le}\,(F_2 \circ F_1)\,H_1\,(G_1 \circ G_2) ,
\]
an entry whose hole names a source proposition. This is where templates pay off. The entries of an object form do not depend on the atom they will be applied to, so object forms compose by syntax. A $\ev{trans}$ between two object coercions is pushed into their entries, which is the cut elimination of Section~\ref{sec:idea} for objects. In the fragment no template proves an equality, so between typed forms the middle telescope has only inclusions and $H_2$ is some $\ev{le}\,j$. If $H_2$ reads an equality, the substitution is undefined.

Between typed forms the last equation is reached only by $\form{top}$ followed by an object form out of $\top$, and the composite $\form{top}$ is typed (Lemma~\ref{lem:mini-combine}).

\begin{figure}[t]
\small
\color{tgtcol}
\[
\begin{array}{@{}lcl@{}}
\mathit{hnf}(\ev{refl}\,T) & = & \form{id} \\
\mathit{hnf}(\ev{eqToLe}\,\varphi) & = & \form{id} \\
\mathit{hnf}(\ev{top}\,T) & = & \form{top} \\
\mathit{hnf}(\ev{bot}\,T) & = & \form{bot} \\
\mathit{hnf}(\ev{pi}\,d\,c) & = & \form{pi}\,d\,c \\
\mathit{hnf}(\ev{obj}\,\tel\,m) & = & \form{obj}\,(\mathit{entries}(m)) \\
\mathit{hnf}(\ev{trans}\,e\,f) & = & \mathit{hnf}(e) \circ \mathit{hnf}(f) \\
\mathit{hnf}(\ev{member}\,a\,e\,i) & = & G \quad \text{where } \tat{\mathit{viewThrough}(\mathit{hnf}(e), a)}{i}{\form{le}\,G} \\[3pt]
\mathit{entries}([\,\tau_0, \ldots, \tau_{n-1}\,]) & = & [\,\mathit{entry}(\tau_0), \ldots, \mathit{entry}(\tau_{n-1})\,] \\
\mathit{entry}(\ev{le}\,q\,H\,q') & = & \form{le}\,(\mathit{side}(q))\,H\,(\mathit{side}(q')) \\
\mathit{side}(\ev{none}) & = & \form{id} \\
\mathit{side}(\ev{some}\,e) & = & \mathit{hnf}(e) \\[3pt]
\mathit{view}(x) & = & [\,\form{eq}, \ldots, \form{eq}\,] \text{ with one } \form{eq} \text{ per witness of } \sigma(x) \\
\mathit{view}(\ev{cast}\,a\,e) & = & \mathit{viewThrough}(\mathit{hnf}(e), a) \\[3pt]
\mathit{viewThrough}(\form{id}, a) & = & \mathit{view}(a) \\
\mathit{viewThrough}(\form{obj}\,\mathit{Es}, a) & = & \mathit{at}_{\mathit{view}(a)}(\mathit{Es}) \\
\mathit{viewThrough}(F, a) & = & \tnil \quad \text{for } F \in \{\form{pi}\,d\,c, \form{top}, \form{bot}\} \\[3pt]
\mathit{at}_V([\,E_0, \ldots, E_{n-1}\,]) & = & [\,\mathit{at}_V(E_0), \ldots, \mathit{at}_V(E_{n-1})\,] \\
\mathit{at}_V(\form{le}\,F\,(\ev{le}\,j)\,G) & = & \form{le}\,((F \circ F') \circ G) \quad \text{if } \tat{V}{j}{\form{le}\,F'} \\
\mathit{at}_V(\form{le}\,F\,H\,G) & = & \form{le}\,(F \circ G) \quad \text{if } H \in \{\ev{eq}\,j, \ev{eqSym}\,j\} \text{ and } \tat{V}{j}{\form{eq}} \\[3pt]
\mathit{chain}(x) & = & \form{id} \\
\mathit{chain}(\ev{cast}\,a\,e) & = & \mathit{chain}(a) \circ \mathit{hnf}(e)
\end{array}
\]
\color{black}
\caption{The normalizer of the fragment, over a fixed store $\sigma$, which the machine writes as a subscript, $\mathit{chain}_\sigma$. Composition $\circ$ is defined in the text. Through $\form{top}$ and $\form{pi}$ the view is empty, since $\top$ and a function type have no propositions, and typed atoms never reach $\form{bot}$. Equality evidence needs no form, because its endpoints are equal after following definitions (Theorem~\ref{thm:mini-canon}, item~2).}
\label{fig:mini-normalizer}
\end{figure}

Figure~\ref{fig:mini-normalizer} gives the normalizer, a family of mutually recursive functions over a fixed store. Each is partial: it fails when a lookup finds the wrong kind of thing, which Section~\ref{sec:mini-theorem} shows does not happen for typed evidence. In the fragment the recursion is structural. On $\ev{member}\,a\,e\,i$ the normalizer computes the view of the atom $a$ through the head form of $e$ and returns its $i$-th entry.

\begin{example}[A selection through a literal]\label{ex:mini-sel}
Let the store hold $\sigma(t) = \tgt{\nu(\self.\,A = \mathit{Int})}$, so $\Gamma(t) = \mathit{Lit}$. The atom $a = \ev{cast}\,t\,\mathit{litCo}$ has the type $R$ of the declaration $\dtyp{A}{\mathit{Int}}{\mathit{Int}}$, and
\[
\ev{member}\ a\ (\ev{refl}\,R)\ 1 \ :\ \sel{t}{A} \le \mathit{Int}
\]
reads its upper bound. It is $\mathit{hi}$ of Section~\ref{sec:mini-evidence} after allocation has replaced $t$ by $a$. The normalizer computes the view of $a$ through $\form{id}$, which is the view of $t$ through the head form of $\mathit{litCo}$. That head form is
\[
\mathit{hnf}(\mathit{litCo}) = \form{obj}\,[\,\form{le}\,\form{id}\,(\ev{eqSym}\,0)\,\form{id},\ \form{le}\,\form{id}\,(\ev{eq}\,0)\,\form{id}\,] .
\]
The view of $t$ is $[\,\form{eq}\,]$. Applying the entries fills both holes with $\form{id}$, so the view of $a$ is $[\,\form{le}\,\form{id},\ \form{le}\,\form{id}\,]$. The evidence normalizes to $\form{id}$: the selection $\sel{t}{A}$ is its definition.
\end{example}

\begin{example}[The argument of $\mathit{id}$]\label{ex:mini-view}
The argument of $\mathit{id}$ is $b = \ev{cast}\,a\,\mathit{inst}$, with $a$ the atom of Example~\ref{ex:mini-sel}. The head form of $\mathit{inst}$ keeps its sides as forms:
\[
\mathit{hnf}(\mathit{inst}) = \form{obj}\,[\,\form{le}\,\form{bot}\,(\ev{le}\,0)\,\form{id},\ \form{le}\,\form{id}\,(\ev{le}\,1)\,\form{top}\,] .
\]
The view of $b$ pushes the view $[\,\form{le}\,\form{id},\ \form{le}\,\form{id}\,]$ of $a$ through it. Entry~0 gives $\form{le}\,((\form{bot} \circ \form{id}) \circ \form{id})$ and entry~1 gives $\form{le}\,((\form{id} \circ \form{id}) \circ \form{top})$, so $\mathit{view}(b) = [\,\form{le}\,\form{bot},\ \form{le}\,\form{top}\,]$. This is all that the body of $\mathit{id}$ could learn about its argument: $\bot \le \sel{t}{A} \le \top$, and nothing about $\mathit{Int}$.

The machine never computes the chain of $b$, since $b$ is never applied. The chain still shows how two object forms compose, and Example~\ref{ex:mini-bad} compares it with the bad-bounds case. It is $\mathit{hnf}(\mathit{litCo}) \circ \mathit{hnf}(\mathit{inst})$, once the $\form{id}$ of $t$ drops out. Entry~0 of $\mathit{hnf}(\mathit{inst})$ names position~0, where $\mathit{hnf}(\mathit{litCo})$ has $\form{le}\,\form{id}\,(\ev{eqSym}\,0)\,\form{id}$, and entry~1 names position~1. Substitution gives
\[
\mathit{chain}(b) = \form{obj}\,[\,\form{le}\,\form{bot}\,(\ev{eqSym}\,0)\,\form{id},\ \form{le}\,\form{id}\,(\ev{eq}\,0)\,\form{top}\,] .
\]
Its holes name the equality of the literal, and the sides of $\mathit{inst}$ sit at its ends. Applied to $\mathit{view}(t) = [\,\form{eq}\,]$, it gives $\mathit{view}(b)$ again.
\end{example}

\FloatBarrier
\subsection{Typed forms}\label{sec:mini-typed}

\begin{figure}[t]
\small
\color{tgtcol}
\begin{mathpar}
\inferrule{\resolve{\Gamma}{S} = \bot}{\formj{\Gamma}{\form{bot}}{S}{T}}
\and
\inferrule{\resolve{\Gamma}{T} = \top}{\formj{\Gamma}{\form{top}}{S}{T}}
\and
\inferrule{\resolve{\Gamma}{S} = \resolve{\Gamma}{T}}{\formj{\Gamma}{\form{id}}{S}{T}}
\and
\inferrule{\resolve{\Gamma}{S} = \pit{S_1}{T_1} \\ \resolve{\Gamma}{T} = \pit{S_2}{T_2} \\ \lej{\Gamma}{d}{S_2}{S_1} \\ \lej{\Gamma, x : S_2}{c}{T_1}{T_2}}
  {\formj{\Gamma}{\form{pi}\,d\,c}{S}{T}}
\and
\inferrule{\resolve{\Gamma}{S} = \obj{\tel_1} \\ \resolve{\Gamma}{T} = \obj{\tel_2} \\ \entj{\Gamma}{\mathit{Es}}{\tel_1}{\tel_2}}
  {\formj{\Gamma}{\form{obj}\,\mathit{Es}}{S}{T}}
\\
\inferrule{ }{\entj{\Gamma}{\tnil}{\tel_1}{\tnil}}
\and
\inferrule{\entj{\Gamma}{\mathit{Es}}{\tel_1}{\tel_2} \\ \holej{\tel_1}{H}{X}{Y} \\ \Gamma \models F : S \le X \text{ (side)} \\ \Gamma \models G : Y \le T \text{ (side)}}
  {\entj{\Gamma}{\tcons{\mathit{Es}}{\form{le}\,F\,H\,G}}{\tel_1}{\tcons{\tel_2}{\ple{S}{T}}}}
\\
\inferrule{ }{\Gamma \models \form{id} : X \le X \text{ (side)}}
\and
\inferrule{\formj{\Gamma}{F}{U_1}{U_2}}{\Gamma \models F : \weak{U_1} \le \weak{U_2} \text{ (side)}}
\\
\inferrule{ }{\viewjr{\Gamma}{r}{\tnil}{\tnil}}
\and
\inferrule{\viewjr{\Gamma}{r}{V}{\tel} \\ \formj{\Gamma}{F}{\inst{S}{r}}{\inst{T}{r}}}
  {\viewjr{\Gamma}{r}{\tcons{V}{\form{le}\,F}}{\tcons{\tel}{\ple{S}{T}}}}
\and
\inferrule{\viewjr{\Gamma}{r}{V}{\tel} \\ \resolve{\Gamma}{\inst{S}{r}} = \resolve{\Gamma}{\inst{T}{r}}}
  {\viewjr{\Gamma}{r}{\tcons{V}{\form{eq}}}{\tcons{\tel}{\peq{S}{T}}}}
\end{mathpar}
\color{black}
\caption{Typed forms, typed entries, typed sides, and typed views in the fragment. A view is typed against a telescope instantiated at a root $r$.}
\label{fig:mini-formtyped}
\end{figure}

Theorem~\ref{thm:mini-canon} has to say what the normalizer returns, in a form that progress can invert by inspection. That is what typed forms are for. Their rules read the endpoints after following definitions, by resolution (Section~\ref{sec:mini-terms}). Example~\ref{ex:mini-sel} shows why: there the evidence for $\sel{t}{A} \le \mathit{Int}$ normalized to $\form{id}$, so $\form{id}$ must count as typed between $\sel{t}{A}$ and $\mathit{Int}$.

\paragraph{Typed forms.}
Figure~\ref{fig:mini-formtyped} gives the typing of forms, written $\models$ to keep it apart from the typing $\vdash$ of evidence. A form is typed for $S \le T$ when its pieces of evidence are typed and the resolved endpoints have the shapes the form promises. Typed entries mirror the typing of templates in Figure~\ref{fig:mini-templates}, with forms on the sides. The first rule for sides types an absent side, $\form{id}$ between syntactically equal types, which may mention the self. The second types a present side, a form between self-free types. So a side with an endpoint that mentions the self is $\form{id}$ between equal types. On its own, $\form{id}$ relates any two types that resolve equally.

Compare the inversion lemma of Section~\ref{sec:textbook}, which Lemma~\ref{lem:mini-shapes} below replaces. Its \rn{Trans} case needed an induction, and in DOT its selection cases would need a context without bad bounds. The typing of forms has neither. It has no transitivity rule, since composition is a lemma about forms (Lemma~\ref{lem:mini-combine}). It has no rule that reads a declaration, since endpoints are compared by resolution, which reads definitions.

\begin{lemma}[Shapes]\label{lem:mini-shapes}
If $\formj{\Gamma}{F}{S}{T}$ then $S$ resolves to $\bot$, or $T$ resolves to $\top$, or $S$ and $T$ resolve equally, or both resolve to function types, or both resolve to object types. In particular there is no typed form for $\top \le \bot$, none from an object type to a function type, and none from a function type to an object type other than $\top$.
\end{lemma}

The proof is a look at Figure~\ref{fig:mini-formtyped}. For $\top \le \bot$: $\top$ is $\obj{\tnil}$, which is neither $\bot$ nor a name, so it resolves to itself, and $\bot$ resolves to $\bot$. $\form{bot}$ would need $\top$ to resolve to $\bot$, and $\form{top}$ would need $\bot$ to resolve to $\top$. $\form{id}$ would need them to be equal, and $\form{pi}$ and $\form{obj}$ would need both to be function types or both object types.

In the full calculus some object types and function types are related (Section~\ref{sec:mini-omitted}), so only the first consequence carries over (Corollary~\ref{cor:shapes}).

\FloatBarrier
\subsection{The normalization theorem}\label{sec:mini-theorem}

We call evidence closed when it is typed in the context of a typed store, so that every free variable is a location. The theorem needs the store. Over the opaque parameter $x : D$ of the overview, $\ev{member}\,x\,(\ev{refl}\,D)\,0 : \top \le \sel{x}{A}$ is typed. But $\sel{x}{A}$ resolves to itself, and Lemma~\ref{lem:mini-shapes} leaves no form from $\top$ to it.

\begin{theorem}[Normalization of closed evidence]\label{thm:mini-canon}
Let $\storej{\sigma}{\Gamma}$.
\begin{enumerate}
\item If $\lej{\Gamma}{e}{S}{T}$ then $\mathit{hnf}(e)$ is defined and $\formj{\Gamma}{\mathit{hnf}(e)}{S}{T}$.
\item If $\eqj{\Gamma}{\varphi}{S}{T}$ then $\resolve{\Gamma}{S} = \resolve{\Gamma}{T}$.
\item If $\morj{\Gamma}{m}{\tel}{\tel'}$ then $\mathit{entries}(m)$ is defined and typed from $\tel$ to $\tel'$.
\item If $\atomj{\Gamma}{a}{S}$ then $\mathit{view}(a)$ is defined, $\resolve{\Gamma}{S} \neq \bot$, and $\mathit{view}(a)$ is typed at $\tel$ at the root of $a$ whenever $\resolve{\Gamma}{S} = \obj{\tel}$.
\end{enumerate}
\end{theorem}

The proof is a mutual induction on the typing derivations, one item per judgment. Three cases need lemmas about forms. $\ev{trans}$ composes two forms (Lemma~\ref{lem:mini-combine}), and $\ev{member}$ and a cast atom push a view through an object form (Lemma~\ref{lem:mini-apply}). Item~4 says that whatever object type an atom has, the value at its root realizes it, and the view says how. It also carries an invariant with no textbook counterpart, that the type of an atom never resolves to $\bot$, which the $\ev{member}$ case needs.

\begin{lemma}[Composition]\label{lem:mini-combine}
If $\formj{\Gamma}{F}{S}{M}$ and $\formj{\Gamma}{G}{M}{T}$, then $F \circ G$ is defined and $\formj{\Gamma}{F \circ G}{S}{T}$.
\end{lemma}

\begin{proof}[Proof sketch]
By induction on the sum of the sizes of $F$ and $G$, by cases on the first equation that applies.
\begin{itemize}
\item $\form{id} \circ G = G$. The form $\form{id}$ says that $S$ and $M$ resolve equally. Typedness reads endpoints only through resolution, so $G$ is typed from $S$ as well. The case $F \circ \form{id}$ is symmetric.
\item $\form{bot} \circ G = \form{bot}$ and $F \circ \form{top} = \form{top}$ keep the endpoint that the form constrains.
\item Two function forms, with $S$, $M$, and $T$ resolving to $\pit{S_1}{T_1}$, $\pit{S_2}{T_2}$, and $\pit{S_3}{T_3}$. The domain evidence composes to $\ev{trans}\,d_2\,d_1 : S_3 \le S_1$. The codomain evidence $c_1$ is typed under a parameter of type $S_2$. Substituting the cast atom $\ev{cast}\,x\,\weak{d_2}$ for that parameter types it under a parameter of type $S_3$. This uses a substitution lemma for evidence, proven before normalization. The substitution does not change any type, because a cast atom has the same root as the variable it replaces, and types only see roots.
\item Two object forms, $\form{obj}\,\mathit{Es}_1 \circ \form{obj}\,\mathit{Es}_2$. The middle telescope is the same for both. It is the target of the typed form $\form{obj}\,\mathit{Es}_1$, so all its propositions are inclusions. Take an entry $\form{le}\,F_2\,H_2\,G_2$ of $\mathit{Es}_2$, whose hole names the inclusion $j$ of the middle telescope. The typed entry of $\mathit{Es}_1$ at $j$ proves that inclusion from the source as $\form{le}\,F_1\,H_1\,G_1$. The composite entry $\form{le}\,(F_2 \circ F_1)\,H_1\,(G_1 \circ G_2)$ is typed by the induction hypothesis on the sides. A present side is self-free, so composing two present sides happens between types outside the object. An absent side composed with another side gives that side.
\item $\form{top} \circ \form{obj}\,\mathit{Es}$, with the object form out of $\top$. It has no entries, since $\top$ has no propositions for a hole to name. So its target resolves to $\top$ as well, and the composite $\form{top}$ is typed.
\item Every other pair has incompatible shapes at $M$. For example, $\form{pi}\,d\,c \circ \form{obj}\,\mathit{Es}$ would need $M$ to resolve both to a function type and to an object type. \qedhere
\end{itemize}
\end{proof}

\begin{lemma}[Application to a view]\label{lem:mini-apply}
If $\entj{\Gamma}{\mathit{Es}}{\tel_1}{\tel_2}$ and $\viewjr{\Gamma}{r}{V}{\tel_1}$, then $\mathit{at}_V(\mathit{Es})$ is defined and is a view typed at $\tel_2$ at the root $r$.
\end{lemma}

\begin{proof}[Proof sketch]
Entry by entry. Take $\form{le}\,F\,H\,G$ for the target inclusion $\ple{S}{T}$, where $H$ reads $\ple{X}{Y}$ in $\tel_1$. The view has at the position of $H$ either a form $F'$ typed from $\inst{X}{r}$ to $\inst{Y}{r}$, or $\form{eq}$, and then $\form{id}$ is typed there. A present side is self-free, so instantiating it at $r$ changes nothing. An absent side stays the identity at $r$. Two uses of Lemma~\ref{lem:mini-combine} type $(F \circ F') \circ G$ from $\inst{S}{r}$ to $\inst{T}{r}$.
\end{proof}

\begin{proof}[Proof of Theorem~\ref{thm:mini-canon}]
By mutual induction on the typing derivations.
\begin{itemize}
\item $\ev{member}\,a\,e\,i$, with $a : S_a$ and $e : S_a \le \obj{\tel}$. Item 4 gives a view of $a$, and item 1 a typed form $F$ for $e$. The form $F$ is $\form{id}$ or $\form{obj}$. It is not $\form{bot}$, because $S_a$ does not resolve to $\bot$ by item 4. It is not $\form{top}$, because $\tel$ has an $i$-th proposition and is not empty. It is not $\form{pi}$, because $\obj{\tel}$ is not a function type. If $F = \form{id}$, the view of $a$ is typed at $\tel$ by item 4. If $F = \form{obj}\,\mathit{Es}$, it is typed at $\tel$ by Lemma~\ref{lem:mini-apply}. Its $i$-th entry is a form typed from $\inst{S'}{\rt{a}}$ to $\inst{T'}{\rt{a}}$, which is the conclusion of the rule.
\item A variable $x$: the store holds a literal or a lambda at $x$, and $\Gamma(x)$ is its type. That is an object type or a function type, so it does not resolve to $\bot$. For a literal, the view lists $\form{eq}$ for each equality $\peq{\sel{\self}{A}}{W(A)}$, and it is typed because $\sel{x}{A}$ resolves like its definition.
\item A cast atom $\ev{cast}\,a\,e$ with $e : S_a \le T$: item 1 gives a typed form $F$ for $e$. By item 4, $S_a$ does not resolve to $\bot$. So $F$ is not $\form{bot}$, and $T$ does not resolve to $\bot$ either, since a form into such a type is $\form{bot}$ or $\form{id}$. Through $\form{id}$ or $\form{obj}$ the view is typed as in the case of $\ev{member}$. Through $\form{top}$ it is empty, which is typed at the empty telescope of $\top$. Through $\form{pi}$ there is nothing to type, since a function type has no telescope.
\item $\ev{refl}$, $\ev{top}$, $\ev{bot}$: the form is typed directly. $\ev{eqToLe}\,\varphi$ gives $\form{id}$, typed by item 2. $\ev{pi}\,d\,c$ gives $\form{pi}\,d\,c$, whose pieces are the typed premises.
\item $\ev{obj}\,\tel\,m$: item 3 types the entries. A present side is self-free evidence typed in $\Gamma$, so item 1 applies to it.
\item $\ev{trans}\,e\,f$: item 1 for both premises and Lemma~\ref{lem:mini-combine}.
\item $\ev{def}\,x\,A$: resolution of a defined name is resolution of its definition. The other equality rules are immediate. \qedhere
\end{itemize}
\end{proof}

\paragraph{Why the induction is structural.}
Section~\ref{sec:idea} compared normalization to cut elimination, with transitivity as the cut. In cut elimination, removing a cut can create new cuts that are not subterms of the proof. Normalization proofs for typed $\lambda$-calculi need a logical relation for the same reason. Neither happens here. A function form keeps its domain and codomain as evidence, so composing two function forms builds new evidence without normalizing it. The machine normalizes it later, after an application has replaced the parameter by an atom over the store. A template reads the self only through its hole, so composing two object forms substitutes entries into entries and never meets an elimination. Section~\ref{sec:why-structural} describes an earlier design where it did.

Theorem~\ref{thm:mini-canon} is about one piece of evidence, but the machine normalizes the chain of casts on the atom it applies. So progress needs the theorem in that form.

\begin{theorem}[Chain of casts]\label{thm:mini-chain}
If $\storej{\sigma}{\Gamma}$ and $\atomj{\Gamma}{a}{S}$, then $\mathit{chain}(a)$ is defined and $\formj{\Gamma}{\mathit{chain}(a)}{\Gamma(\rt{a})}{S}$.
\end{theorem}

The proof is an induction on the atom. A variable has the chain $\form{id}$ at its own type, and a cast composes the chain underneath with the form of its evidence, which Theorem~\ref{thm:mini-canon} types, by Lemma~\ref{lem:mini-combine}.

\begin{corollary}[Consistency and canonical forms]\label{cor:mini}
Let $\storej{\sigma}{\Gamma}$.
\begin{enumerate}
\item There is no $e$ with $\lej{\Gamma}{e}{\top}{\bot}$.
\item If $\atomj{\Gamma}{a}{\pit{S}{T}}$, then $\sigma(\rt{a})$ is a lambda, and $\mathit{chain}(a)$ is $\form{id}$ or a function form.
\end{enumerate}
\end{corollary}

\begin{proof}
The first item is Theorem~\ref{thm:mini-canon} and Lemma~\ref{lem:mini-shapes}. For the second, Theorem~\ref{thm:mini-chain} gives a form typed from $\Gamma(\rt{a})$ to $\pit{S}{T}$. The type of a location is the type of its stored value, a function type for a lambda and an object type for a literal. By Lemma~\ref{lem:mini-shapes} there is no form from an object type to a function type, so the stored value is a lambda. Its type is a function type, and the forms between two function types are $\form{id}$ and $\form{pi}$, since $\form{bot}$, $\form{top}$, and $\form{obj}$ need an endpoint of another shape.
\end{proof}

Item~2 is the canonical forms lemma of Section~\ref{sec:textbook}, over a store, and it is all that progress needs. It rests on there being no closed evidence from an object type to a function type, which does the job of the consistency lemma of System FC (Section~\ref{sec:mini-overview}). Item~1 is not needed for safety. It is the bad-bounds argument of the overview in general form, and Example~\ref{ex:mini-bad} traces it.

\begin{example}[Bad bounds, revisited]\label{ex:mini-bad}
Here is the overview again, with every step justified. Let $\lambda(x : D).\,u$ be applied to an atom $a$ of type $D$ with root $z$. By Theorem~\ref{thm:mini-chain} the chain of $a$ is a typed form from $\Gamma(z)$ to $D$. It is the head form of the composite coercion $c$ of the overview. Lemma~\ref{lem:mini-shapes} has no form from a function type to $D$, so $z$ holds a literal $\tgt{\nu(\self.\,W)}$, and $\Gamma(z)$ is its precise type $\obj{\tgt{\mathit{tel}(W)}}$, which has only equalities. $D$ has inclusions, so the two endpoints are distinct object types, and $D$ is not $\top$. So Figure~\ref{fig:mini-formtyped} leaves only an object form, with one entry per proposition of $D$.

Entry~0 proves $\ple{\top}{\sel{\self}{A}}$ in three links, $\top \le X \sqsubseteq Y \le \sel{\self}{A}$. A side with an endpoint that mentions the self is absent and links that type only to itself, and a present side is self-free. So $Y$ is $\sel{\self}{A}$, and $X$ is self-free, as $\top$ is. The hole reads an equality $\peq{\sel{\self}{B}}{W(B)}$ of the literal, forwards or backwards. The name $\sel{\self}{B}$ mentions the self, so it is $Y$, and the hole reads backwards. Hence $B$ is $A$, the literal defines $A$, and $X$ is its witness $T = W(A)$, which is self-free. The translation produces only such literals (Section~\ref{sec:mini-translation}), and for any other literal entry~0 cannot exist. The left side of entry~0 is a form $F_0$ typed from $\top$ to $T$, and its right side is $\form{id}$. Entry~1 is the mirror image, with a right side $G_1$ typed from $T$ to $\bot$. With $j$ the position of $A$ in $W$, the chain is
\[
\form{obj}\,[\,\form{le}\,F_0\,(\ev{eqSym}\,j)\,\form{id},\ \ \form{le}\,\form{id}\,(\ev{eq}\,j)\,G_1\,] \qquad \text{with } \formj{\Gamma}{F_0}{\top}{T} \text{ and } \formj{\Gamma}{G_1}{T}{\bot} .
\]
By Lemma~\ref{lem:mini-combine}, $F_0 \circ G_1$ is typed from $\top$ to $\bot$, which Lemma~\ref{lem:mini-shapes} excludes. So no such $a$ exists, and the lambda is never applied in a typed state.

Compare the chain of $b$ in Example~\ref{ex:mini-view}, $\form{obj}\,[\,\form{le}\,\form{bot}\,(\ev{eqSym}\,0)\,\form{id},\ \form{le}\,\form{id}\,(\ev{eq}\,0)\,\form{top}\,]$ from $\mathit{Lit}$ to $\mathit{Abs}$. It has the same shape. In both, the definition sits between the bounds, $\bot \le \mathit{Int} \le \top$ for good bounds and $\top \le T \le \bot$ for bad ones. Only the second composes to an inclusion that has no form.
\end{example}

\FloatBarrier
\subsection{Preservation and progress}\label{sec:mini-safety}

With the corollary, preservation and progress have the textbook shape.

\begin{theorem}[Preservation]
If a typed state steps, the new state is typed at the same type, weakened under the new location when the step allocates.
\end{theorem}

The cases that are not about application use weakening, substitution of an atom for a variable, and the typing of stripped casts. Allocation also uses transparency: a term typed with a binder opaque stays typed with it transparent. The running example of Section~\ref{sec:mini-machine} showed why no type changes, and Section~\ref{sec:metatheory} gives the case in full.

The application through a cast atom with chain $\form{pi}\,d\,c$ is the case worth reading. The store holds $\lambda(z : S_0).\,t_0$ at the root, so the root has the type $\pit[z]{S_0}{T_0}$, where $T_0$ is the type of the body. By Theorem~\ref{thm:mini-chain} the chain is a form typed from $\pit[z]{S_0}{T_0}$ to the type $\pit[z]{S}{T}$ of $a$. So $d : S \le S_0$, and $c : T_0 \le T$ under a parameter of type $S$. The argument atom $b$ of the rule has type $S$, so $\ev{cast}\,b\,d$ has type $S_0$ and can be substituted into the body, which gives a term of type $\inst{T_0}{\rt{b}}$. The cast $\inst{c}{b}$ takes that type to $\inst{T}{\rt{b}}$, which is the type of the application. With chain $\form{id}$ the two function types resolve equally, and since neither is a name they are equal, so the body can be used as it is.

In the running example, $\mathit{id}$ returns $\lambda(y : \sel{t}{A}).\,y$, and a let stores it at $g$. The program of Section~\ref{sec:mini-evidence} applies the bare $g$, so its application is $\beta$. A derivation that retypes $\src{g}$ by \rn{All} instead ends in $(\ev{cast}\,g\,(\ev{pi}\,\mathit{lo}\,\weak{\mathit{hi}}))\,n$, with the codomain evidence weakened under the parameter. Write $\mathit{lo}'$ and $\mathit{hi}'$ for $\mathit{lo}$ and $\mathit{hi}$ after allocation has replaced $t$ by $\ev{cast}\,t\,\mathit{litCo}$. The chain is $\form{pi}\,\mathit{lo}'\,\weak{\mathit{hi}'}$, and the state steps to $\ev{cast}\,(\ev{cast}\,n\,\mathit{lo}')\,\mathit{hi}'$. The domain evidence casts $n$ up to the parameter type $\sel{t}{A}$, and the codomain evidence casts the result back to $\mathit{Int}$. The program of Section~\ref{sec:mini-evidence} reaches the same atom by $\beta$.

\begin{theorem}[Progress]
A typed state is final or steps.
\end{theorem}

Every rule of the machine applies to its form of state without a condition, except application. Take an application through an atom $a$ of function type. Corollary~\ref{cor:mini} says that the store holds a lambda at the root and that the chain of casts is $\form{id}$ or $\form{pi}\,d\,c$. One of the three application rules applies.

In preservation the textbook proof inverts $\src{S_1 \to T_1 \sub S_2 \to T_2}$ to obtain the subtypings of the domains and the codomains. Here they are the components $d$ and $c$ of the function form, and the machine inserts them as casts. No derivation of evidence typing is ever inverted.

\FloatBarrier
\subsection{Translating the fragment}\label{sec:mini-translation}

So far \FCdot is a calculus of its own. To say anything about DOT, every typed program of the source fragment must become a typed \FCdot program that erases to it. The source fragment has the types
\[
\src{S, T ::= \top \mid \bot \mid \dtyp{A}{S}{T} \mid x.A \mid \dall{x}{S}{T}} .
\]
We write $\trans{\cdot}$ around a source phrase for its translation, and the postfix $\llbracket\cdot\rrbracket$ stays instantiation. Types translate homomorphically, except that a declaration becomes an object type:
\[
\begin{array}{@{}rcl@{}}
\trans{\src{\top}} & = & \tgt{\top} \\
\trans{\src{\bot}} & = & \tgt{\bot} \\
\trans{\dsel{x}{A}} & = & \sel{x}{A} \\
\trans{\dall{x}{S}{T}} & = & \pit{\trans{S}}{\trans{T}} \\
\trans{\dtyp{A}{S}{T}} & = & \obj{[\,\ple{\weak{\trans{S}}}{\sel{\self}{A}},\ \ple{\sel{\self}{A}}{\weak{\trans{T}}}\,]}
\end{array}
\]
The source fragment keeps the subtyping rules \rn{Top}, \rn{Bot}, \rn{Refl}, \rn{Trans}, \rn{Typ}, \rn{Sel-<:}, \rn{<:-Sel}, and \rn{All}, and the typing rules \rn{Var}, \rn{All-I}, \rn{All-E}, \rn{Let}, and \rn{Sub}. In place of \rn{\{\}-I} it has one rule for literals, $\dtyj{\Gamma}{\dnew{x}{\ddeftyp{A}{T}}}{\dtyp{A}{T}{T}}$, where $\src{T}$ does not mention $\src{x}$. The mechanized source has no such rule, so this translation is not mechanized as it stands. A subtyping derivation becomes evidence, and a typing of a variable becomes an atom rooted at that variable. \rn{Top}, \rn{Bot}, \rn{Refl}, and \rn{Trans} map to $\ev{top}$, $\ev{bot}$, $\ev{refl}$, and $\ev{trans}$. \rn{All} maps to $\ev{pi}$, \rn{Sub} to a cast, and \rn{All-E} to the application of atoms.

Two rules read declarations. If $h$ derives $\src{x : \dtyp{A}{S}{T}}$ and $\trans{h}_a$ is its atom, the two selection rules eliminate at that atom:
\[
\begin{array}{@{}rcl@{}}
\trans{\rn{<:-Sel}\,h} & = & \ev{member}\ \trans{h}_a\ (\ev{refl}\,\trans{\dtyp{A}{S}{T}})\ 0 \ :\ \trans{S} \le \sel{x}{A} \\
\trans{\rn{Sel-<:}\,h} & = & \ev{member}\ \trans{h}_a\ (\ev{refl}\,\trans{\dtyp{A}{S}{T}})\ 1 \ :\ \sel{x}{A} \le \trans{T}
\end{array}
\]
The rule \rn{Typ}, which relates two declarations of $A$, becomes the object coercion of Section~\ref{sec:mini-evidence} with the translated premises as sides:
\[
\trans{\rn{Typ}\,d_1\,d_2} = \ev{obj}\,(\mathit{tel}\,\dtyp{A}{S_1}{T_1})\,[\,\ev{le}\,(\ev{some}\,\trans{d_1})\,(\ev{le}\,0)\,\ev{none},\ \ev{le}\,\ev{none}\,(\ev{le}\,1)\,(\ev{some}\,\trans{d_2})\,] ,
\]
where $\mathit{tel}\,\src{T}$ is the telescope of $\trans{T}$. A literal $\dnew{x}{\ddeftyp{A}{T}}$ typed at $\dtyp{A}{T}{T}$ becomes the literal $\tgt{\nu(\self.\,A = \trans{T})}$ under the cast $\mathit{litCo}_{\trans{T}}$.

The program of Section~\ref{sec:dot}, with the literal typed at $\dtyp{A}{\mathit{Int}}{\mathit{Int}}$, translates to the program of Section~\ref{sec:mini-evidence}. In it, $\mathit{lo}$ and $\mathit{hi}$ are the translations of the selection rules at $t$, and $\mathit{inst}$ is the translation of \rn{Typ} with \rn{Bot} and \rn{Top}. Section~\ref{sec:dot} typed $\src{g\,n}$ by raising $\src{n}$ to $\src{t.A}$ and lowering the result to $\mathit{Int}$. The derivation that retypes $\src{g}$ by \rn{All} instead gives the program of Section~\ref{sec:mini-safety}. The translation follows the derivation, and both programs erase to the same term.

A bad-bounds program, $\src{\lambda(x : \dtyp{A}{\top}{\bot}).\ \klet y = x \kin y}$, whose body is typed at $\dtyp{B}{\mathit{Int}}{\mathit{Int}}$, retypes $x$ by
\[
\src{\dtyp{A}{\top}{\bot} \sub \top \sub x.A \sub \bot \sub \dtyp{B}{\mathit{Int}}{\mathit{Int}}} .
\]
It translates to $\lambda(x : D).\ \klet y = \ev{cast}\,x\,e \kin y$ with
\[
\tgt{e} = \ev{trans}\ (\ev{top}\,D)\ (\ev{trans}\ (\ev{member}\,x\,(\ev{refl}\,D)\,0)\ (\ev{trans}\ (\ev{member}\,x\,(\ev{refl}\,D)\,1)\ (\ev{bot}\,R_B))) ,
\]
where $R_B = \trans{\dtyp{B}{\mathit{Int}}{\mathit{Int}}}$. The two middle steps are the eliminations of the overview, which read the two propositions of the type of $x$. It goes through $\top \le \sel{x}{A} \le \bot$ by ordinary $\ev{trans}$ steps, and the type checker never looks for bad bounds. Section~\ref{sec:translation} gives the whole translation.

Safety of the source fragment follows through erasure. A typed source program translates to a typed \FCdot program with the same erasure. The two run side by side, and preservation and progress of the target show that the source never gets stuck. Section~\ref{sec:safety} gives this argument for full DOT.

\FloatBarrier
\subsection{What the miniature leaves out}\label{sec:mini-omitted}

The full calculus adds the following to the fragment, and Sections~\ref{sec:target} and~\ref{sec:canonical} explain each.
\begin{itemize}
\item Fields add term members, with a proposition that says a field is present, evidence that proves it, and a projection that reads it. The type of a field of $x$ is the block name of $x$ at the field's label.
\item Recursive types let the bounds mention the self, as in $\dmu{x}{\dtyp{A}{\bot}{x.B}}$, and add atoms that fold and unfold an object type at its root. These atoms change the type of an atom without evidence, so the chain of casts is typed with object types opened at the root.
\item Intersections concatenate telescopes and add a pairing of two coercions out of one type and of two typings of one variable.
\item A type that is not a declaration, such as $\src{x.A}$ or a function type, enters a telescope as one self-bound, which says that the object itself is included in the type. So an object type with a self-bound may be included in a function type, and a function type in an object type whose propositions are all self-bounds. Self-bounds are also why the full calculus composes $\form{top}$ with an object form differently.
\item Two more kinds of template, for equalities and for chaining two templates, and the elimination of equalities by $\ev{member}$. The translation of DOT uses none of them.
\end{itemize}
None of these changes what preservation and progress ask of the theorem: a shape.
